\documentclass[10pt,a4paper]{amsart}
\usepackage{amsmath,amssymb,amsthm,graphicx}
\usepackage{verbatim}
\newenvironment{reference}{\comment}{\endcomment}
%\newenvironment{reference}{}{}
\newenvironment{slogan}{\comment}{\endcomment}
\newenvironment{history}{\comment}{\endcomment}

% For commutative diagrams we use Xy-pic
\usepackage[all]{xy}

% We use 2cell for 2-commutative diagrams.
\xyoption{2cell}
% Theorem environments.
%
\theoremstyle{plain}
\newtheorem{theorem}[subsection]{Theorem}
\newtheorem{proposition}[subsection]{Proposition}
\newtheorem{lemma}[subsection]{Lemma}

\theoremstyle{definition}
\newtheorem{definition}[subsection]{Definition}
\newtheorem{example}[subsection]{Example}
\newtheorem{exercise}[subsection]{Exercise}
\newtheorem{situation}[subsection]{Situation}

\theoremstyle{remark}
\newtheorem{remark}[subsection]{Remark}
\newtheorem{remarks}[subsection]{Remarks}

\numberwithin{equation}{subsection}

% Macros
%
\def\lim{\mathop{\mathrm{lim}}\nolimits}
\def\colim{\mathop{\mathrm{colim}}\nolimits}
\def\Spec{\mathop{\mathrm{Spec}}}
\def\Hom{\mathop{\mathrm{Hom}}\nolimits}
\def\Ext{\mathop{\mathrm{Ext}}\nolimits}
\def\SheafHom{\mathop{\mathcal{H}\!\mathit{om}}\nolimits}
\def\SheafExt{\mathop{\mathcal{E}\!\mathit{xt}}\nolimits}
\def\Sch{\mathit{Sch}}
\def\Mor{\mathop{\mathrm{Mor}}\nolimits}
\def\Ob{\mathop{\mathrm{Ob}}\nolimits}
\def\Sh{\mathop{\mathit{Sh}}\nolimits}
\def\NL{\mathop{N\!L}\nolimits}
\def\CH{\mathop{\mathrm{CH}}\nolimits}
\def\proetale{{pro\text{-}\acute{e}tale}}
\def\etale{{\acute{e}tale}}
\def\QCoh{\mathit{QCoh}}
\def\Ker{\mathop{\mathrm{Ker}}}
\def\Im{\mathop{\mathrm{Im}}}
\def\Coker{\mathop{\mathrm{Coker}}}
\def\Coim{\mathop{\mathrm{Coim}}}

% Boxtimes
%
\DeclareMathSymbol{\boxtimes}{\mathbin}{AMSa}{"02}

%
% Macros for moduli stacks/spaces
%
\def\QCohstack{\mathcal{QC}\!\mathit{oh}}
\def\Cohstack{\mathcal{C}\!\mathit{oh}}
\def\Spacesstack{\mathcal{S}\!\mathit{paces}}
\def\Quotfunctor{\mathrm{Quot}}
\def\Hilbfunctor{\mathrm{Hilb}}
\def\Curvesstack{\mathcal{C}\!\mathit{urves}}
\def\Polarizedstack{\mathcal{P}\!\mathit{olarized}}
\def\Complexesstack{\mathcal{C}\!\mathit{omplexes}}
% \Pic is the operator that assigns to X its picard group, usage \Pic(X)
% \Picardstack_{X/B} denotes the Picard stack of X over B
% \Picardfunctor_{X/B} denotes the Picard functor of X over B
\def\Pic{\mathop{\mathrm{Pic}}\nolimits}
\def\Picardstack{\mathcal{P}\!\mathit{ic}}
\def\Picardfunctor{\mathrm{Pic}}
\def\Deformationcategory{\mathcal{D}\!\mathit{ef}}
\usepackage{cleveref}
\crefname{originallemma}{original lemma}{original lemmas}
\Crefname{originallemma}{Original lemma}{Original lemmas}
\makeatletter
\newlabel{derived-lemma-leray-acyclicity}{{015E}{source}{}{sourcepointer-1}{}}
\newlabel{derived-lemma-leray-acyclicity@cref}{{[originallemma][0][]015E}{[1][0][]source}{}{}{}}
\makeatother
\crefname{originaltheorem}{original theorem}{original theorems}
\Crefname{originaltheorem}{Original theorem}{Original theorems}
\makeatletter
\newlabel{injectives-theorem-K-injective-embedding-grothendieck}{{079P}{source}{}{sourcepointer-2}{}}
\newlabel{injectives-theorem-K-injective-embedding-grothendieck@cref}{{[originaltheorem][0][]079P}{[1][0][]source}{}{}{}}
\makeatother
\crefname{originalreference}{original reference}{original references}
\Crefname{originalreference}{Original reference}{Original references}
\makeatletter
\newlabel{lemma-injective-pushforward}{{09TL}{source}{}{sourcepointer-3}{}}
\newlabel{lemma-injective-pushforward@cref}{{[originalreference][0][]09TL}{[1][0][]source}{}{}{}}
\makeatother
\crefname{originalreference}{original reference}{original references}
\Crefname{originalreference}{Original reference}{Original references}
\makeatletter
\newlabel{lemma-quasi-coherence-direct-image}{{08D5}{source}{}{sourcepointer-4}{}}
\newlabel{lemma-quasi-coherence-direct-image@cref}{{[originalreference][0][]08D5}{[1][0][]source}{}{}{}}
\makeatother
\crefname{originallemma}{original lemma}{original lemmas}
\Crefname{originallemma}{Original lemma}{Original lemmas}
\makeatletter
\newlabel{spaces-morphisms-lemma-pushforward}{{03M9}{source}{}{sourcepointer-5}{}}
\newlabel{spaces-morphisms-lemma-pushforward@cref}{{[originallemma][0][]03M9}{[1][0][]source}{}{}{}}
\makeatother
\crefname{originallemma}{original lemma}{original lemmas}
\Crefname{originallemma}{Original lemma}{Original lemmas}
\makeatletter
\newlabel{spaces-morphisms-lemma-quasi-compact-quasi-separated-permanence}{{073B}{source}{}{sourcepointer-6}{}}
\newlabel{spaces-morphisms-lemma-quasi-compact-quasi-separated-permanence@cref}{{[originallemma][0][]073B}{[1][0][]source}{}{}{}}
\makeatother
\crefname{originalproposition}{original proposition}{original propositions}
\Crefname{originalproposition}{Original proposition}{Original propositions}
\makeatletter
\newlabel{spaces-properties-proposition-coherator}{{077V}{source}{}{sourcepointer-7}{}}
\newlabel{spaces-properties-proposition-coherator@cref}{{[originalproposition][0][]077V}{[1][0][]source}{}{}{}}
\makeatother
\title{Derived pushforward of quasi-coherent sheaves computed inside the quasi-coherent category}
\author{The Stacks Project Authors}
\date{stacks-project:spaces-perfect.tex:local-snapshot-20261012; GNU FDL 1.2 or later}
\begin{document}
\maketitle
\noindent\textit{Editorial note.} All delivered mathematical text is the author's own, copied from the pinned Stacks Project chapter file spaces-perfect.tex; no statement or proof step was rewritten, reordered or removed.\par\medskip
\noindent\textit{Editorial note.} The counted proof is the author's complete argument: the construction and the comparison are both fully written out; the bounded-below case is discharged by Leray acyclicity with the injective pushforward vanishing\par\medskip
\noindent\textit{Editorial note.} The label lemma-Noetherian-pushforward names the counted statement, and the proof environment immediately following it in the source is delivered verbatim as its proof.\par\medskip
\noindent\textit{Editorial note.} The counted statement is carried by the author's own lemma environment in the chapter source.\par\medskip
\noindent\textit{Allowed standard prerequisites (editorial).}\begin{itemize}
\item derived pushforward Rf\_* on modules over algebraic spaces and its right derived functor viewpoint
\item Grothendieck abelian categories admit K-injective resolutions by complexes of injectives
\item the uniform vanishing bound N of the quasi-coherence direct image lemma
\item truncation short exact sequences of complexes
\item Leray acyclicity for bounded below complexes of f\_*-acyclic objects and the injective-pushforward lemma
\end{itemize}\medskip
\ifdefined\hypertarget\hypertarget{sourcepointer-1}{}\fi
\noindent\textit{Source reference.} Original lemma 015E: sources/stacks-chow/tags.txt, source lines 1505--1505. Prerequisite pointer invoked at proof line(s) 3030; it names a standard fact of a sibling chapter that the counted argument uses, and every step of the counted proof built on it is retained verbatim in the delivered body.\par\medskip
\ifdefined\hypertarget\hypertarget{sourcepointer-2}{}\fi
\noindent\textit{Source reference.} Original theorem 079P: sources/stacks-chow/tags.txt, source lines 9279--9279. Prerequisite pointer invoked at proof line(s) 2999; it names a standard fact of a sibling chapter that the counted argument uses, and every step of the counted proof built on it is retained verbatim in the delivered body.\par\medskip
\ifdefined\hypertarget\hypertarget{sourcepointer-3}{}\fi
\noindent\textit{Source reference.} Original reference 09TL: sources/stacks-spaces-perfect/pinned-original-source.tex, source lines 2917--2917. Prerequisite pointer invoked at proof line(s) 3031; it names a standard fact of the same chapter that the counted argument uses, and every step of the counted proof built on it is retained verbatim in the delivered body.\par\medskip
\ifdefined\hypertarget\hypertarget{sourcepointer-4}{}\fi
\noindent\textit{Source reference.} Original reference 08D5: sources/stacks-spaces-perfect/pinned-original-source.tex, source lines 790--790. Prerequisite pointer invoked at proof line(s) 3013; it names a standard fact of the same chapter that the counted argument uses, and every step of the counted proof built on it is retained verbatim in the delivered body.\par\medskip
\ifdefined\hypertarget\hypertarget{sourcepointer-5}{}\fi
\noindent\textit{Source reference.} Original lemma 03M9: sources/stacks-chow/tags.txt, source lines 4741--4741. Prerequisite pointer invoked at proof line(s) 2990; it names a standard fact of a sibling chapter that the counted argument uses, and every step of the counted proof built on it is retained verbatim in the delivered body.\par\medskip
\ifdefined\hypertarget\hypertarget{sourcepointer-6}{}\fi
\noindent\textit{Source reference.} Original lemma 073B: sources/stacks-chow/tags.txt, source lines 9055--9055. Prerequisite pointer invoked at proof line(s) 2988; it names a standard fact of a sibling chapter that the counted argument uses, and every step of the counted proof built on it is retained verbatim in the delivered body.\par\medskip
\ifdefined\hypertarget\hypertarget{sourcepointer-7}{}\fi
\noindent\textit{Source reference.} Original proposition 077V: sources/stacks-chow/tags.txt, source lines 9215--9215. Prerequisite pointer invoked at proof line(s) 2994; it names a standard fact of a sibling chapter that the counted argument uses, and every step of the counted proof built on it is retained verbatim in the delivered body.\par\medskip
\begin{lemma}
\label{lemma-Noetherian-pushforward}
Let $S$ be a scheme.
Let $f : X \to Y$ be a morphism of Noetherian algebraic spaces over $S$.
Then $f_*$ on quasi-coherent sheaves has a right derived
extension
$\Phi : D(\QCoh(\mathcal{O}_X)) \to D(\QCoh(\mathcal{O}_Y))$
such that the diagram
$$
\xymatrix{
D(\QCoh(\mathcal{O}_X)) \ar[d]_{\Phi} \ar[r] &
D_\QCoh(\mathcal{O}_X) \ar[d]^{Rf_*} \\
D(\QCoh(\mathcal{O}_Y)) \ar[r] &
D_\QCoh(\mathcal{O}_Y)
}
$$
commutes.
\end{lemma}

\begin{proof}
Since $X$ and $Y$ are Noetherian the morphism is quasi-compact
and quasi-separated (see
Morphisms of Spaces, Lemma
\ref{spaces-morphisms-lemma-quasi-compact-quasi-separated-permanence}).
Thus $f_*$ preserve quasi-coherence, see
Morphisms of Spaces, Lemma \ref{spaces-morphisms-lemma-pushforward}.
Next, let $K$ be an object of $D(\QCoh(\mathcal{O}_X))$.
Since $\QCoh(\mathcal{O}_X)$ is a Grothendieck abelian category
(Properties of Spaces, Proposition
\ref{spaces-properties-proposition-coherator}), we can
represent $K$ by a K-injective complex $\mathcal{I}^\bullet$
such that each $\mathcal{I}^n$ is an injective object of
$\QCoh(\mathcal{O}_X)$, see
Injectives, Theorem
\ref{injectives-theorem-K-injective-embedding-grothendieck}.
Thus we see that the functor $\Phi$ is defined by setting
$$
\Phi(K) = f_*\mathcal{I}^\bullet
$$
where the right hand side is viewed as an object of
$D(\QCoh(\mathcal{O}_Y))$. To finish the proof of the lemma
it suffices to show that the canonical map
$$
f_*\mathcal{I}^\bullet \longrightarrow Rf_*\mathcal{I}^\bullet
$$
is an isomorphism in $D(\mathcal{O}_Y)$. To see this it suffices to
prove the map induces an isomorphism on cohomology sheaves. Pick any
$m \in \mathbf{Z}$. Let $N = N(X, Y, f)$ be as in
Lemma \ref{lemma-quasi-coherence-direct-image}.
Consider the short exact sequence
$$
0 \to \sigma_{\geq m - N - 1}\mathcal{I}^\bullet \to
\mathcal{I}^\bullet \to \sigma_{\leq m - N - 2}\mathcal{I}^\bullet \to 0
$$
of complexes of quasi-coherent sheaves on $X$. By
Lemma \ref{lemma-quasi-coherence-direct-image}
we see that the cohomology sheaves of
$Rf_*\sigma_{\leq m - N - 2}\mathcal{I}^\bullet$
are zero in degrees $\geq m - 1$. Thus we see that
$R^mf_*\mathcal{I}^\bullet$ is isomorphic to
$R^mf_*\sigma_{\geq m - N - 1}\mathcal{I}^\bullet$.
In other words, we may assume that $\mathcal{I}^\bullet$
is a bounded below complex of injective objects of
$\QCoh(\mathcal{O}_X)$.
This case follows from Leray's acyclicity lemma
(Derived Categories, Lemma \ref{derived-lemma-leray-acyclicity})
with required vanishing because of Lemma \ref{lemma-injective-pushforward}.
\end{proof}

\end{document}
